\documentclass[11pt]{article}
\usepackage{amsmath,amssymb}
\usepackage[margin=1in]{geometry}
\usepackage{hyperref}

\title{Physics of elkpy's Fortran-patch extensions}
\author{}
\date{}

\begin{document}
\maketitle

\tableofcontents

\part{Per-species scaling of the spin-orbit coupling term}
\label{part:soc}

\section{Setting: second-variational spin-orbit coupling in Elk}

Elk solves the Kohn-Sham problem in a second-variational scheme: a scalar-relativistic
Hamiltonian is diagonalized first, and spin-dependent terms (spin polarization,
spin-orbit coupling) are added afterward in the basis of scalar-relativistic
eigenstates. When \texttt{spinorb = .true.}, the spin-orbit term added inside each
muffin-tin sphere is the Koelling-Harmon (1977) approximation to the spin-orbit part of
the Dirac equation \cite{koelling1977}:
\begin{equation}
  \hat H_{\rm soc}(r) = f_{\rm soc}(r)\,\hat{\mathbf L}\cdot\boldsymbol\sigma,
  \qquad
  f_{\rm soc}(r) = \frac{1}{\bigl(2 M(r) c\bigr)^2}\,\frac{1}{r}\,
    \frac{\partial V_s(r)}{\partial r},
  \label{eq:hsoc}
\end{equation}
with the relativistic mass correction
\begin{equation}
  M(r) = 1 + \frac{1}{2c^2}\bigl(E - V_s(r)\bigr)\Big|_{E=0},
\end{equation}
where $V_s(r)$ is the spherical part of the Kohn-Sham effective potential in that
muffin-tin, $c$ is the speed of light (\texttt{solsc} in Elk's atomic-unit convention),
$\hat{\mathbf L}$ is the orbital angular momentum operator, and $\boldsymbol\sigma$ are
the Pauli matrices. This is exactly what Elk's \texttt{gensocfr.f90} computes per atom
(verified against \texttt{vendor/elk/src/gensocfr.f90}, whose own docstring cites the
same Koelling-Harmon reference).

Elk exposes one global scale on top of Eq.~\eqref{eq:hsoc}, the input parameter
\texttt{socscf} (default $1.0$): the radial coefficient actually evaluated is
\begin{equation}
  c_{\rm so} = \frac{y_{00}\cdot \texttt{socscf}}{4\,c^2},
  \qquad
  f_{\rm soc}(r) = \frac{c_{\rm so}}{r\,M(r)^2}\,\frac{\partial V_s(r)}{\partial r},
  \label{eq:socscf}
\end{equation}
the same scalar multiplying $\hat H_{\rm soc}$ for every atom in the cell, regardless
of species. Elk's manual states this knob's purpose plainly: it is used ``to enhance
the effect of spin-orbit coupling in order to accurately determine the magnetic
anisotropy energy (MAE)'' — i.e.\ it is a phenomenological correction factor, fitted
per material to compensate for whatever the scalar-relativistic/second-variational
approximation to Eq.~\eqref{eq:hsoc} misses relative to a full four-component Dirac
treatment or experiment, not a quantity derived from first principles.

\section{Why a single global scale is the wrong shape for multi-species cells}

$f_{\rm soc}(r)$ in Eq.~\eqref{eq:hsoc} is dominated by the region near the nucleus,
where $V_s(r)$ is steep; the resulting spin-orbit splitting grows strongly with atomic
number $Z$ (the same near-nuclear physics behind atomic fine-structure splitting,
heuristically $\sim Z^4$ for hydrogenic orbitals). A \texttt{socscf} value fitted to
correct the SOC-driven MAE of one species in a compound — typically a heavy element
whose scalar-relativistic treatment most understates true spin-orbit effects — is not,
in general, the correction a light species in the same cell needs. Because
\texttt{socscf} is a single number applied uniformly to every atom
(Eq.~\eqref{eq:socscf}), Elk upstream cannot express ``scale species A's spin-orbit
term but leave species B's alone.''

\section{elkpy's generalization}

elkpy promotes the scale in Eq.~\eqref{eq:socscf} from a single global number to a
per-species one, $\texttt{socscf} \to s_{\rm is}$, so that for an atom of species $is$:
\begin{equation}
  c_{\rm so}^{(is)} = \frac{y_{00}\cdot s_{is}}{4\,c^2}, \qquad
  s_{is} = \begin{cases}
    \texttt{soc\_scale[is]} & \text{if overridden}\\
    \texttt{socscf} & \text{otherwise (Elk's existing global default)}
  \end{cases}
  \label{eq:persp}
\end{equation}
Nothing about the physical content of Eq.~\eqref{eq:hsoc} changes — the radial
derivative $\partial V_s/\partial r$, the relativistic mass factor $M(r)$, and the
$\hat{\mathbf L}\cdot\boldsymbol\sigma$ structure are exactly Elk's existing
Koelling-Harmon term. Only the multiplicative prefactor becomes species-resolved, so a
fitted correction can be scoped to the species it was actually fitted for instead of
leaking onto every other species in the cell.

\subsection{Implementation}

\texttt{gensocfr.f90} already loops over every atom (\texttt{do ias=1,natmtot}) to
evaluate $f_{\rm soc}(r)$ for that atom's species, so
\texttt{patches/0001-per-species-soc-scale.patch} only inserts the species-dependent
lookup of Eq.~\eqref{eq:persp} inside that existing loop — it does not restructure the
term itself:
\begin{itemize}
  \item \texttt{modmain.f90}: a new array \texttt{socscfsp(maxspecies)}, sentinel
    $< 0$ meaning ``not overridden by elkpy, fall back to \texttt{socscf}''.
  \item \texttt{readinput.f90}: a new input block \texttt{elkpy\_socscale} (one
    \texttt{is~value} pair per line) that populates \texttt{socscfsp}.
  \item \texttt{gensocfr.f90}: inside the existing per-atom loop, \texttt{scso =
    socscf}, then \texttt{if (socscfsp(is) >= 0) scso = socscfsp(is)}, and $c_{\rm so}$
    is computed from \texttt{scso} instead of \texttt{socscf} directly.
\end{itemize}
This is the smallest possible hook into existing control flow, per the isolation
constraint in \texttt{docs/design.md} §8 — no new files were strictly needed here since
the per-atom loop already existed to hang the lookup on.

\subsection{Verification}

Verified against a real compiled Elk binary (\texttt{tests/test\_calculation\_soc.py}):
\begin{itemize}
  \item Single-species cell: \texttt{soc\_scale} reproduces the result of the
    equivalent global \texttt{socscf} value exactly (Eq.~\eqref{eq:persp} reduces to
    Eq.~\eqref{eq:socscf} when only one species is present).
  \item Two-species cell: scaling one species' \texttt{soc\_scale} changes that
    species' contribution independently of the other, confirming the override is
    applied per-species rather than falling back to a single effective global value.
\end{itemize}

\section{How to use in code}

\begin{verbatim}
from elkpy import Calculation

calc = Calculation(
    structure,
    spinpol=True,
    spinorb=True,
    soc_scale={"Fe": 1.5},   # scale only Fe's spin-orbit term; other species
                             # keep Elk's default global socscf (1.0 unless
                             # overridden separately via extra_blocks)
)
e = calc.get_energy()
\end{verbatim}
\texttt{soc\_scale} requires \texttt{spinorb=True} (elkpy raises otherwise, since a
scale on a disabled term has no physical effect), every key must name a species present
in the \texttt{Structure}, and every value must be $\geq 0$.

\part{Berry curvature via the Wilson-loop (Fukui-Hatsugai-Suzuki) method}

\section{Setting: Berry curvature and the Chern number}

For a family of Bloch states $|n(\mathbf k)\rangle$ (a non-degenerate band, or a
gapped multiplet of bands) over the Brillouin zone, the Berry connection and the
associated field strength (Berry curvature) are
\begin{equation}
  A_\mu(\mathbf k) = \langle n(\mathbf k) | \partial_\mu | n(\mathbf k) \rangle, \qquad
  F_{12}(\mathbf k) = \partial_1 A_2(\mathbf k) - \partial_2 A_1(\mathbf k),
  \label{eq:berry}
\end{equation}
with $\mu=1,2$ two reciprocal-space directions and $\partial_\mu = \partial/\partial
k_\mu$. When the Fermi energy lies in a gap, the (2D) Hall conductance is
$\sigma_{xy} = -(e^2/h)\sum_n c_n$, where the Chern number of band $n$,
\begin{equation}
  c_n = \frac{1}{2\pi i} \int_{T^2} d^2k\, F_{12}(\mathbf k),
  \label{eq:chern}
\end{equation}
is an integer topological invariant (Thouless-Kohmoto-Nightingale-den Nijs,
Phys.\ Rev.\ Lett.\ \textbf{49}, 405 (1982)) -- nonzero only when $A_\mu$ cannot be
chosen as a single smooth gauge over the whole (torus) Brillouin zone, which is exactly
the topological obstruction the Chern number counts. $|n(\mathbf k)\rangle$ carries an
arbitrary $\mathbf k$-dependent phase (the gauge freedom of quantum mechanics);
$A_\mu$ depends on that choice, $F_{12}$ and $c_n$ do not.

\section{The Wilson-loop (Fukui-Hatsugai-Suzuki) discretization}

In a real DFT calculation, wavefunctions exist only on a discrete $\mathbf k$-mesh, and
the phase Elk's (or any) diagonalizer returns at each mesh point is an arbitrary,
numerically-determined gauge choice -- not smooth across the mesh, not known in closed
form. Naively finite-differencing Eq.~\eqref{eq:berry} in that raw gauge does not work.
Fukui, Hatsugai and Suzuki (FHS; \emph{Chern Numbers in Discretized Brillouin Zone:
Efficient Method of Computing (Spin) Hall Conductances}, J.\ Phys.\ Soc.\ Jpn.\
\textbf{74}, 1674 (2005), arXiv:cond-mat/0503172) give a construction that is exactly
gauge invariant on the discretized zone, for {\it any} choice of per-$\mathbf k$-point
phase, by never differentiating a phase directly and instead multiplying wavefunction
overlaps ("Wilson loops") around each mesh plaquette -- the same idea used for Berry
curvature in tight-binding/lattice models, e.g.\ \texttt{pyqula}'s
\texttt{berry\_curvature} (which this implementation was cross-checked against for
convention, per this project's development-practices policy).

For a mesh point $\mathbf k_\ell$ and the two lattice directions $\hat\mu$ ($\mu=1,2$)
spanning the loop, and a (possibly multi-band) manifold $\psi(\mathbf k) =
(|n_1(\mathbf k)\rangle,\dots,|n_M(\mathbf k)\rangle)$, define the U(1) link variable
\begin{equation}
  U_\mu(\mathbf k_\ell) = \frac{\det\, \psi^\dagger(\mathbf k_\ell)\,
    \psi(\mathbf k_\ell+\hat\mu)}{\bigl|\det\, \psi^\dagger(\mathbf k_\ell)\,
    \psi(\mathbf k_\ell+\hat\mu)\bigr|},
  \label{eq:link}
\end{equation}
a single unit-modulus complex number per link, built from the $M \times M$ overlap
matrix between the band manifold at $\mathbf k_\ell$ and at its mesh neighbour
$\mathbf k_\ell+\hat\mu$ (the Abelian, $M=1$, case is Eq.~7 of FHS; Eq.~\eqref{eq:link}
here is their non-Abelian generalization, their Eq.~16). The lattice field strength
around the plaquette is
\begin{equation}
  \tilde F_{12}(\mathbf k_\ell) = \ln\!\bigl[
    U_1(\mathbf k_\ell)\, U_2(\mathbf k_\ell+\hat 1)\,
    U_1(\mathbf k_\ell+\hat 2)^{-1}\, U_2(\mathbf k_\ell)^{-1} \bigr],
  \qquad -\pi < \frac{1}{i}\tilde F_{12}(\mathbf k_\ell) \le \pi,
  \label{eq:fieldstrength}
\end{equation}
i.e.\ the principal-branch phase of the product of the four link variables walked
around the plaquette. Because $\tilde F_{12}$ is built entirely from gauge-covariant
overlaps (never a bare derivative of a phase), it is exactly invariant under
$|n_a(\mathbf k)\rangle \to e^{-i\lambda_a(\mathbf k)}|n_a(\mathbf k)\rangle$ for {\it
any} $\mathbf k$-dependent phases $\lambda_a$ -- this is the property elkpy's
\texttt{tests/test\_berry\_gauge\_invariance.py} checks directly (Section~7.4). The
lattice Chern number is the sum over the mesh,
\begin{equation}
  \tilde c_n = \frac{1}{2\pi i} \sum_\ell \tilde F_{12}(\mathbf k_\ell),
  \label{eq:latticechern}
\end{equation}
and FHS prove it is \emph{exactly} an integer for any mesh spacing (not just in a
continuum limit), converging to Eq.~\eqref{eq:chern} as the mesh is refined, correct
even on a coarsely discretized zone provided the \emph{admissibility condition}
$|\tilde F_{12}(\mathbf k_\ell)| < \pi$ holds at every plaquette (their Eq.~14) -- a
value approaching $\pi$ signals a mesh too coarse to trust.

\section{elkpy's implementation}

\subsection{What Elk already provides, and why it isn't reused directly}

Elk's Wannier90 export (task 550, \texttt{writew90mmn.f90}) computes exactly this kind
of quantity -- an overlap matrix $M_{mn}(\mathbf k,\mathbf k+\mathbf b) = \langle
\psi_m(\mathbf k)|\psi_n(\mathbf k+\mathbf b)\rangle$ between neighbouring $\mathbf
k$-points on a mesh -- via \texttt{genwfsvp} (expand the second-variational
wavefunction, muffin-tin plus interstitial) and \texttt{genolpq} (the overlap itself,
including the $e^{i\mathbf b\cdot\mathbf r}$ phase factor needed inside the muffin-tins
for the cell-periodic part). It cannot be used as-is here: which neighbours
$\mathbf k+\mathbf b$ to use is decided by \texttt{wannier\_setup}, a call into the
external Wannier90 library that is only a stub (\texttt{w90\_stub.f90}, ``libwannier not
or improperly installed'') unless Elk is separately linked against libwannier90 -- not
part of this project's build (\texttt{build-config/make.inc}).

elkpy sidesteps this rather than adding a libwannier90 build dependency: the two
Wilson-loop directions $\hat\mu$ are, by construction, exactly two of the three
\texttt{ngridk} mesh generators, so the neighbour of $\mathbf k_\ell$ in direction $\mu$
is simply $\mathbf k_\ell + \hat\mu$ in lattice coordinates -- $1/\texttt{ngridk}(\mu)$
of a reciprocal lattice vector -- with no neighbour-shell search needed at all.

\subsection{Fortran: exporting the overlap matrices (task 9000)}

\texttt{patches/0002-berry-curvature-wilson-loop.patch} adds one new task, reserved in
the high, clearly-unused range per \texttt{docs/design.md}~\S8, and one new additive
file, \texttt{elkpy\_berry.f90} (\texttt{elkpy\_berrycurv} subroutine):
\begin{itemize}
  \item \texttt{modmain.f90}: four new integers (\texttt{elkpy\_berry\_dir1/dir2},
    the two loop directions, and \texttt{elkpy\_berry\_ist0/ist1}, the requested band
    window).
  \item \texttt{readinput.f90}: a new input block \texttt{elkpy\_berry} populating
    them, with validation (directions distinct and in $\{1,2,3\}$, \texttt{ist0}
    $\ge 1$, \texttt{ist1} $\ge$ \texttt{ist0}).
  \item \texttt{elk.f90}: one line, \texttt{case(9000): call elkpy\_berrycurv}, plus a
    documentation block for the new input block (same pattern as
    \texttt{elkpy\_socscale}, Part~\ref{part:soc}).
  \item \texttt{elkpy\_berry.f90} (new file): for every point of the full,
    non-reduced \texttt{ngridk} mesh and for each of the two requested directions,
    reuses \texttt{genwfsvp}/\texttt{genolpq} exactly as \texttt{writew90mmn.f90} does,
    but against the directly-constructed neighbour $\mathbf k_\ell+\hat\mu$ instead of
    a Wannier90-supplied shell, and writes the resulting $M^{(\mu)}(\mathbf k_\ell)$
    matrices (restricted to the requested band window) plus, for a gap-closing check,
    the eigenvalues of the window and one band on each side of it, to
    \texttt{ELKPY\_BERRY.OUT}.
\end{itemize}
This requires \texttt{reducek=0} (set automatically by
\texttt{Calculation.get\_berry\_curvature()}): with symmetry reduction on, eigenvectors
are only computed/stored on the irreducible wedge of the mesh, not the full mesh the
Wilson loop walks.

No Wilson-loop arithmetic happens in Fortran at all -- Eqs.~\eqref{eq:link}
through~\eqref{eq:latticechern} are evaluated entirely in Python
(\texttt{elkpy.parsers.berry}), so that part is unit-testable against synthetic overlap
matrices without running Elk (Section~7.4).

\subsection{Python: the Wilson loop and Chern number}

\texttt{elkpy.parsers.berry.compute\_berry\_curvature} evaluates
Eqs.~\eqref{eq:link}--\eqref{eq:latticechern} directly with \texttt{numpy}, at every
point of the mesh (looping over the direction not spanned by the loop as independent
2D slices), and additionally reports $\max_\ell |\tilde F_{12}(\mathbf k_\ell)|$, the
admissibility diagnostic. \texttt{check\_gap} verifies, from the exported boundary
eigenvalues, that the requested band window stays separated from the rest of the
spectrum at every mesh point -- required for Eq.~\eqref{eq:link}'s multi-band
determinant to be well defined (FHS's relaxed gap-opening condition for the
non-Abelian case: $E_n(\mathbf k) \ne E_{n'}(\mathbf k)$ for $n$ in the window,
$n'$ outside it, for all $\mathbf k$) -- and raises rather than silently returning a
number computed through a near-degeneracy.

\subsection{Verification}

\begin{itemize}
  \item \emph{Gauge invariance} (\texttt{tests/test\_berry\_gauge\_invariance.py}, no
    Elk needed): applying a random per-$\mathbf k$-point, per-band diagonal unitary
    gauge transformation $M(\mathbf k,\mathbf k') \to D(\mathbf k)^\dagger\, M(\mathbf
    k,\mathbf k')\, D(\mathbf k')$ to synthetic overlap matrices leaves the computed
    flux and Chern number exactly unchanged (to floating-point precision) -- validates
    the arithmetic itself, independent of whether the overlaps are physical. This does
    \emph{not} pin the overall sign of the result: replacing every overlap matrix $M$
    by $M^*$ flips the sign of every $\tilde F_{12}$ but leaves the gauge-cancellation
    algebra (Section~6, the derivation showing the four link phases cancel exactly)
    completely unaffected, so a conjugation-sign bug would pass this test undetected.
  \item \emph{Sign} (also \texttt{tests/test\_berry\_gauge\_invariance.py}): the sign of
    the Python-side arithmetic is instead pinned directly against Eq.~\eqref{eq:fieldstrength}
    with hand-constructed single-band overlap matrices of known phase -- one test checks
    a link entering the numerator role ($U_1(\mathbf k_\ell)$), one checks a link entering
    the (inverted) denominator role ($U_1(\mathbf k_\ell+\hat 2)$), each asserting the
    resulting flux has the correspondingly correct sign. This pins that
    \texttt{compute\_berry\_curvature} implements Eq.~\eqref{eq:fieldstrength} with the
    right overall sign \emph{given} overlap matrices in the
    $M(a,b)=\langle\psi_a(\mathbf k)|\psi_b(\mathbf k+\hat\mu)\rangle$ convention -- it
    does not by itself verify that \texttt{elkpy\_berry.f90} produces matrices in that
    convention.
  \item \emph{Convention (Fortran)}: the overlap convention
    ($M(a,b)=\langle\psi_a(\mathbf k)|\psi_b(\mathbf k+\hat\mu)\rangle$, i.e.\
    \texttt{conjg(oq(b,a))} of \texttt{genolpq}'s own output) was pinned down directly
    from the documented \texttt{zgemv} BLAS semantics of \texttt{genolpq.f90} --
    derived independently twice, agreeing both times -- rather than by an empirical
    runtime check: since Eq.~\eqref{eq:link} only ever uses $\det M$, which is
    transpose-invariant, only the conjugation (not the row/column convention) is
    physically consequential, and no test currently exercises the real
    Fortran-to-Python sign chain end to end. An attempted empirical check (requesting
    the overlap of $\Gamma$ with itself via a full reciprocal lattice vector, on a
    $1\times1\times1$ mesh, which should trivially equal the identity matrix) did not
    hold, traced to an apparent limitation of \texttt{getevecfv}'s G-vector relabelling
    for that specific degenerate case -- not exercised by any real calculation, since
    \texttt{elkpy\_berry.f90} now rejects \texttt{ngridk}$<2$ in either loop direction
    outright (a single mesh point gives a zero-area plaquette regardless). A real
    multi-point mesh ($2\times2\times2$) showed consistent $|\det M|$ across both
    boundary-wrapping and non-wrapping neighbour pairs, evidence against (though not
    proof against) the same issue affecting realistic usage.
  \item \emph{Physics} (\texttt{tests/test\_calculation\_berry.py}, needs a built
    binary): bulk Si's 4 valence bands are a trivial (non-topological) insulator --
    the Chern number on every $k_3$ slice comes out $\sim 10^{-19}$ (floating-point
    zero, not merely ``close to an integer''), with the admissibility diagnostic
    ($\max|\tilde F_{12}|\lesssim 2\times 10^{-3}$ on a $4\times4\times4$ mesh) well
    inside FHS's validity regime. This exercises the full Fortran-to-Python pipeline
    end to end, but -- since Si's Chern number is $0=-0$ regardless of sign -- it
    cannot by itself catch a conjugation-sign error; no topologically nontrivial
    reference case is checked here yet.
\end{itemize}

\section{How to use in code: mesh mode}

\begin{verbatim}
from elkpy import Structure

s = Structure(avec, species)
calc = s.get_calculation(xc="PW", ngridk=(4, 4, 4))

result = calc.get_berry_curvature(1, 4, directions=(1, 2))
result["flux"]          # (n1, n2, n3) array, dimensionless plaquette flux
result["chern_number"]  # (n3,) array, one Chern number per k3 slice
result["max_flux"]      # admissibility diagnostic; keep well under pi
\end{verbatim}
\texttt{ist0}/\texttt{ist1} (here 1, 4) give the contiguous, 1-indexed
second-variational band window (e.g.\ the occupied bands); \texttt{directions}
(here the default, $(1,2)$) picks the two \texttt{ngridk} mesh directions the loop is
built from, leaving the third as the free (slice) direction. The window must stay
gapped from the rest of the spectrum at every mesh $\mathbf k$-point -- checked
automatically, raising \texttt{ValueError} otherwise.

\section{Path mode: Berry curvature at arbitrary k-points}

\subsection{Motivation and construction}

Mesh mode only ever evaluates on a periodic mesh covering the whole Brillouin zone --
the only way to get a genuine Chern number (Eq.~\eqref{eq:latticechern} is a sum over a
closed cover of $T^2$), but every Wilson-loop corner then has to be an actual mesh
point, since \texttt{genwfsvp} only knows how to expand a wavefunction from an
eigenvector \emph{already computed and stored} for a specific $k$-point
(\texttt{getevecfv}/\texttt{getevecsv} read from \texttt{EVECFV.OUT}/\texttt{EVECSV.OUT}).
Evaluating curvature along an arbitrary band-structure-style path (e.g.\
$\Gamma$-K-M-K$'$-$\Gamma$) then means either interpolating a mesh or aligning
\texttt{ngridk} to every point of interest by hand.

Path mode instead evaluates one small, independent Wilson loop directly at each
requested point $\mathbf k_0$ -- the same construction \texttt{pyqula}'s own
\texttt{berry\_curvature(h, k, dk=...)} uses for a tight-binding Hamiltonian. Four
corners
\begin{equation}
  \mathbf k_0 - \hat 1\,\texttt{dk} - \hat 2\,\texttt{dk}, \quad
  \mathbf k_0 + \hat 1\,\texttt{dk} - \hat 2\,\texttt{dk}, \quad
  \mathbf k_0 + \hat 1\,\texttt{dk} + \hat 2\,\texttt{dk}, \quad
  \mathbf k_0 - \hat 1\,\texttt{dk} + \hat 2\,\texttt{dk}
  \label{eq:corners}
\end{equation}
are walked cyclically, and the plaquette flux is the phase of the product of the four
edge link variables,
\begin{equation}
  \theta(\mathbf k_0) = \arg\bigl[U_{1\to2}\,U_{2\to3}\,U_{3\to4}\,U_{4\to1}\bigr],
  \qquad U_{i\to j} = \frac{\det M_{ij}}{|\det M_{ij}|}, \quad
  (M_{ij})_{ab} = \langle\psi_a(\text{corner } i)|\psi_b(\text{corner } j)\rangle,
  \label{eq:pathflux}
\end{equation}
and the curvature is $\theta(\mathbf k_0)$ divided by the loop's actual Cartesian area,
$|(2\,\texttt{dk}\,\mathbf b_1)\times(2\,\texttt{dk}\,\mathbf b_2)|$ (not a bare
$\texttt{dk}^2$, which silently assumes an orthogonal reciprocal lattice).
Eq.~\eqref{eq:pathflux} is algebraically identical to Eq.~\eqref{eq:fieldstrength} with
$k_\ell = $ corner 1: since $M_{ji} = M_{ij}^\dagger$ (a direct consequence of
$\langle\psi_a(j)|\psi_b(i)\rangle = \overline{\langle\psi_b(i)|\psi_a(j)\rangle}$), each
``backward'' edge's link variable is exactly the complex conjugate -- i.e.\ the
inverse, since $|U|=1$ -- of the corresponding forward one, so the plain 4-term product
in Eq.~\eqref{eq:pathflux} reduces exactly to FHS's
$U_1(k_\ell)U_2(k_\ell+\hat1)U_1(k_\ell+\hat2)^{-1}U_2(k_\ell)^{-1}$ once corners 2 and 4
are identified with $k_\ell+\hat1$ and $k_\ell+\hat2$. No Chern number is defined for a
single loop (that needs the closed cover), and no automatic gap check runs in this mode
(no eigenvalues are exported, unlike mesh mode's boundary-window eigenvalues) --
degeneracy shows up instead as a singular or ill-conditioned $M_{ij}$, or as a value that
grows rather than converges as \texttt{dk} shrinks (Section~9.4).

\subsection{Implementation: fresh diagonalisation, not a mesh lookup}

The one piece of machinery this needed that mesh mode didn't: \texttt{elkpy\_wfcorner}
(\texttt{elkpy\_berry.f90}) expands a wavefunction at an arbitrary $\mathbf k$ by a
\emph{fresh diagonalisation} of the converged Kohn-Sham Hamiltonian
(\texttt{eveqnfv}/\texttt{eveqnsv}) rather than reading a previously stored eigenvector
-- the same on-the-fly pattern \texttt{src/bandstr.f90} already uses for an arbitrary
band-structure path (\texttt{readstate} $\to$ \texttt{genvsig} $\to$ \texttt{linengy}
$\to$ \texttt{genapwlofr} $\to$ \texttt{gensocfr}, then diagonalise), reusing
\texttt{genwfsvp}'s own on-the-fly $\mathbf G+\mathbf k$-vector setup
(\texttt{gengkvec}/\texttt{gensfacgp}/\texttt{match}) rather than the mesh's persistent
bookkeeping arrays. \texttt{elkpy\_berrycurv\_path} (task 9001) builds the four corners
per requested point, calls \texttt{elkpy\_wfcorner} at each, reuses \texttt{genolpq} for
the four cyclic edges exactly as task 9000 does, and writes the edge overlap matrices to
\texttt{ELKPY\_BERRY\_PATH.OUT}; \texttt{parsers.berry.compute\_berry\_curvature\_path}
evaluates Eq.~\eqref{eq:pathflux} in Python, for the same testability reason as mesh
mode. Because each corner is a one-off diagonalisation from the converged
density/potential, the ground state's own \texttt{ngridk} is irrelevant here -- unlike
mesh mode, this task does not require \texttt{reducek}$=0$, and a $\Gamma$-only ground
state is not a stopgap but simply sufficient, since task 9001 never touches whatever
mesh the ground state happened to converge on.

\subsection{The dk floor}

\texttt{dk} is the accuracy knob, and unlike \texttt{pyqula}'s tight-binding version it
has a floor: each corner is an independent LAPW diagonalisation, and the overlap between
two corners separated by a very small \texttt{dk} is dominated by basis-truncation noise
once \texttt{dk} is small enough -- the plaquette phase becomes a difference of
nearly-identical quantities. There is no universally correct default; the working
procedure is to evaluate one point of interest at a few \texttt{dk} values and look for
where the resulting curvature plateaus (Section~9.4) before trusting a full path --
\texttt{get\_berry\_curvature\_path} does not do this automatically.

\subsection{Verification}

\begin{itemize}
  \item \emph{Gauge invariance and sign}, on synthetic single-loop data, mirroring
    Section~7.4's mesh-mode tests: a random per-corner gauge transform leaves flux and
    curvature exactly unchanged, and the sign of Eq.~\eqref{eq:pathflux} is pinned by
    giving exactly one of the four edges a known phase and checking the flux equals it
    (every edge enters the plain product with the same sign, unlike mesh mode's mixed
    numerator/denominator form -- a difference confirmed by the derivation above, not
    just asserted). The area normalisation is checked against a deliberately
    non-orthogonal (hexagonal-like) \texttt{bvec}, so a bug that only appears for a
    skewed reciprocal lattice isn't masked by testing with an identity metric
    (\texttt{tests/test\_berry\_gauge\_invariance.py}).
  \item \emph{Cross-check against mesh mode}, against a real compiled binary: for
    \texttt{ngridk}$=(2,2,2)$, task 9000's plaquette anchored at $\Gamma$ visits
    $\Gamma,\Gamma+\hat1,\Gamma+\hat1+\hat2,\Gamma+\hat2$ in that cyclic order
    ($\hat\mu=1/\texttt{ngridk}(\mu)$); a path loop centred at
    $\mathbf k_0=(1/4,1/4,0)$ with $\texttt{dk}=1/4$ visits the identical four points in
    the identical order (Eq.~\eqref{eq:corners} with $\mathbf k_0-\hat\mu\,\texttt{dk}$
    etc.\ working out to exactly those corners). On bulk Si, the two modes agree to
    ordinary numerical-noise tolerance ($\sim10^{-5}$, not bit-for-bit, since one path
    reads a converged mesh eigenvector and the other diagonalises fresh) --
    (\texttt{tests/test\_calculation\_berry.py::test\_path\_and\_mesh\_conventions\_agree}).
    This is the end-to-end Fortran-to-Python sign-chain check Section~7.4 notes was
    still missing for mesh mode, though it only catches a discrepancy \emph{between} the
    two modes, not a sign error shared by both.
  \item \emph{Physics}, against a real compiled binary, on monolayer $h$-BN (broken
    sublattice inversion symmetry between B and N, unlike Si -- so K and K$'$ are
    physically inequivalent valleys rather than related by a crystal symmetry): the
    occupied manifold's curvature vanishes at $\Gamma$ and M and is exactly
    antisymmetric between K and K$'$ ($\Omega(\mathrm K)=8.568$,
    $\Omega(\mathrm K')=-8.568$ Bohr$^{-2}$, agreeing to 0.01\%) -- the sign flip
    time-reversal symmetry requires for a non-magnetic crystal, $\Omega(-\mathbf
    k)=-\Omega(\mathbf k)$, and a substantially sharper discriminating check than Si's
    Chern number, which is $0=-0$ regardless of sign. Getting to that result surfaced
    two real usage pitfalls: (1) the occupied-band count must come from Elk's own
    \texttt{EIGVAL.OUT} occupation numbers, not be assumed from a total (core $+$
    valence) electron count -- core electrons are not among the valence bands
    \texttt{nstsv} indexes at all, and getting this wrong silently selects the wrong
    band without any error; (2) a single band's curvature can diverge approaching a
    point where it is degenerate with a \emph{neighbouring occupied} band, not just the
    first unoccupied one, even when the requested window's own outer boundary is safely
    gapped -- $h$-BN's bands 3 and 4 are degenerate exactly at $\Gamma$, so band 4 alone
    diverges there while the full occupied window (bands 1--4 together) does not, since
    the non-Abelian construction (Eq.~\eqref{eq:link}) only requires the window gapped
    from \emph{outside} itself, not internally non-degenerate.
\end{itemize}

\subsection{How to use in code}

\begin{verbatim}
result = calc.get_berry_curvature_path(
    [(0.0, 0.0, 0.0), (1/3, 1/3, 0.0)],  # Gamma, K -- fractional coordinates
    1, 4, directions=(1, 2), dk=0.005,
)
result[0]["curvature"]  # Bohr^-2, one dict per requested k-point
\end{verbatim}
\texttt{kpoints} need not lie on any mesh; \texttt{ist0}/\texttt{ist1}/\texttt{directions}
match mesh mode's meaning. Unlike mesh mode, the requested window is \emph{not}
automatically checked for a closing gap -- see Section~9.4's second pitfall above.

In place of \texttt{kpoints}, a symbolic path can be passed as
\texttt{kpath="GKMG"} (pyqula/ASE style), discretized into \texttt{npoints} points along
the path; each returned point then also carries a \texttt{"distance"} entry (cumulative
Cartesian distance along the path) for plotting against the same $x$-axis convention as
\texttt{get\_bands()}. Unlike \texttt{get\_bands()}/\texttt{get\_phonon\_dispersion()},
whose \texttt{kpath} must be a single connected segment (they hand vertices to Elk's own
\texttt{plot1d} task, which interpolates one continuous line), a disconnected \texttt{','}
path (e.g. \texttt{"GKM,K'G"}, to continue past M into a specific $K'$ zone image rather
than the nearest image a \texttt{plot1d}-style interpolation would pick) is fully
supported here, since task 9001 evaluates every point's Wilson loop independently via
fresh diagonalisation -- there is no interpolation across the break to get wrong:

\begin{verbatim}
result = calc.get_berry_curvature_path(
    ist0=1, ist1=4, directions=(1, 2), dk=0.005,
    kpath="GKM,K'G", npoints=200,
)
result[-1]["distance"]  # cumulative distance along the path, for plotting
\end{verbatim}

\part{Eigenstates and wavefunction overlaps}
\label{part:eigenstates}

\section{Setting: the LAPW/APW+lo generalized eigenvalue problem}

Elk represents each first-variational (scalar-relativistic) Bloch state at a
$\mathbf k$-point in an augmented-plane-wave-plus-local-orbital (APW+lo) basis
$\{\phi_i(\mathbf k,\mathbf r)\}$ that is \emph{not} orthogonal: the basis functions overlap
inside the muffin-tin spheres by construction, so $S_{ij}(\mathbf k) = \langle
\phi_i(\mathbf k)|\phi_j(\mathbf k)\rangle \neq \delta_{ij}$. Diagonalizing the Hamiltonian
therefore means solving the \emph{generalized} eigenvalue problem
\begin{equation}
  H(\mathbf k)\,\mathbf c_n = E_n\, S(\mathbf k)\,\mathbf c_n,
  \label{eq:genevp}
\end{equation}
which is exactly what Elk's \texttt{eveqnfv.f90} does (a generalized Hermitian eigensolver).
The solver normalizes eigenvectors so that $\mathbf c_m^\dagger S(\mathbf k)\, \mathbf c_n =
\delta_{mn}$ -- i.e.\ the physical states $|\psi_m^{\rm fv}(\mathbf k)\rangle = \sum_i
c_{im}\,\phi_i(\mathbf k,\mathbf r)$ \emph{are} orthonormal in the physical sense required of
quantum-mechanical eigenstates, but a \emph{raw} dot product of two coefficient vectors,
$\mathbf c_m^\dagger \mathbf c_n$, is \emph{not} $\delta_{mn}$ in general, since it omits
$S(\mathbf k)$. This is Elk's \texttt{evecfv} -- coefficients in a basis whose own overlap
matrix a caller does not have direct access to, so \texttt{evecfv} alone is not usable for
computing overlaps between states.

\section{Second-variational states: orthonormal, but only within one diagonalization}

The second-variational step builds the spin-dependent Hamiltonian (spin polarization,
spin-orbit coupling, Part~\ref{part:soc}) in the basis of the already-$S$-orthonormal
first-variational eigenstates, $\{|\psi_m^{\rm fv}(\mathbf k)\rangle\}\otimes\{|{\uparrow}
\rangle,|{\downarrow}\rangle\}$ -- itself an ordinary orthonormal basis, with no
non-orthogonal-basis subtlety at this stage (\texttt{eveqnsv.f90} is a standard Hermitian
eigenproblem, not a generalized one). Its eigenvectors, \texttt{evecsv}, therefore do satisfy
plain unitarity,
\begin{equation}
  \texttt{evecsv}^\dagger\,\texttt{evecsv} = \mathbb{1},
  \label{eq:evecsv-unitary}
\end{equation}
verified directly against a real compiled binary
(\texttt{tests/test\_calculation\_eigenstates.py::test\_evecsv\_is\_orthonormal}).

This orthonormality is basis-relative, though, and the basis itself -- the first-variational
eigenstates $\{|\psi_m^{\rm fv}(\mathbf k)\rangle\}$ -- is different at every $\mathbf k$
(different $\mathbf G+\mathbf k$ vectors enter its construction in the first place). So the
\texttt{evecsv} columns returned by two different eigenstate queries -- whether at two
different $\mathbf k$-points, or two independent diagonalizations at literally the same
$\mathbf k$ -- are coefficients in two \emph{different} underlying bases; a raw dot product
between them carries no physical meaning. \texttt{get\_eigenstates()} is documented
accordingly: valid for inspecting a single query's own result (energies, spin/orbital
character via \texttt{evecsv}), not for comparing eigenstates across queries.

\section{Wavefunction overlaps across k-points}

The overlap that \emph{is} meaningful regardless of whether the two $\mathbf k$-points
coincide is the direct overlap integral of the full Bloch wavefunctions themselves,
\begin{equation}
  O_{ab}(\mathbf k_a,\mathbf k_b) = \langle \psi_a(\mathbf k_a) | \psi_b(\mathbf k_b) \rangle
  = \int_{\rm cell} \psi_a^*(\mathbf k_a,\mathbf r)\,\psi_b(\mathbf k_b,\mathbf r)\,d^3r,
  \label{eq:overlap-integral}
\end{equation}
exactly the same quantity as the mesh-neighbour overlap matrix underlying the link variable
of Eq.~\eqref{eq:link} in Part~\ref{part:berry}, generalized from a fixed mesh-step
neighbour $\mathbf k+\hat\mu$ (task 9000) or a small loop's four fixed corners (task 9001) to
an \emph{arbitrary} explicit pair $(\mathbf k_a,\mathbf k_b)$. Computationally it is the same
two-step construction used there: expand each side's second-variational wavefunction into its
muffin-tin-plus-interstitial real-space representation (\texttt{genwfsvp}, or here
\texttt{elkpy\_wfcorner}'s fresh on-the-fly version of it, since $\mathbf k_a$/$\mathbf k_b$
need not lie on any mesh or previously-diagonalized point), then evaluate
Eq.~\eqref{eq:overlap-integral} via \texttt{genolpq}, which folds in the cell-periodic phase
factor $e^{i\mathbf q\cdot\mathbf r}$ ($\mathbf q = \mathbf k_b-\mathbf k_a$) needed inside the
muffin-tins.

\section{elkpy's implementation: an interactive session (task 9002)}

Unlike every other elkpy task, \texttt{get\_eigenstates()}/\texttt{get\_overlap()} are backed
by an \emph{interactive} task, \texttt{elkpy\_eigenstate\_session} (task 9002,
\texttt{patches/0003-eigenstate-session.patch}, new file \texttt{elkpy\_eigenstates.f90}):
after the usual ground-state-resume setup, it loops reading one query per line from standard
input -- \texttt{EIGENSTATES $k_1\,k_2\,k_3$} (Eq.~\eqref{eq:genevp}--\eqref{eq:evecsv-unitary}:
full \texttt{evalsv}/\texttt{evecsv} at an arbitrary point) or \texttt{OVERLAP $k_{1a}\,k_{2a}\,
k_{3a}\,k_{1b}\,k_{2b}\,k_{3b}\,ist_0\,ist_1$} (Eq.~\eqref{eq:overlap-integral} for a
contiguous band window) -- writing each response back to standard output, until a
\texttt{QUIT} command. Both queries reuse \texttt{elkpy\_wfcorner} (Part~\ref{part:berry}),
extended here with optional \texttt{evalsv\_out}/\texttt{evecsv\_out} arguments so an
\texttt{EIGENSTATES} query can retrieve the diagonalization's own eigenvalues/eigenvectors
directly, not only its real-space expansion.

This design -- one long-lived process kept alive across many queries, rather than one process
per query like every other elkpy task -- is an engineering choice, not a physics one; the
rationale (why a persistent worker process rather than an f2py in-memory bridge) is covered in
\texttt{docs/design.md}~\S14, not repeated here.

\subsection{Verified}

Against a real compiled binary (bulk Si): Eq.~\eqref{eq:evecsv-unitary} holds at an arbitrary
$\mathbf k$-point; $O_{ab}(\mathbf k,\mathbf k)$ is the identity matrix, to a tolerance
($\sim 10^{-3}$) set by the \texttt{genwfsvp}/\texttt{genolpq} real-space expansion's own
inherent truncation error (angular-momentum and interstitial $\mathbf G$-vector cutoffs), not
machine precision -- observed at this order even for a single, non-degenerate band at a
generic $\mathbf k$-point, so it is a property of the numerical scheme, not a bug; a session
issuing the same queries twice returns identical results (the session actually stays usable
across many queries); and, cross-checking two independent Fortran code paths the same way
Part~\ref{part:berry}'s path/mesh agreement test does, $O_{ab}(\Gamma,\Gamma+\hat e_1)$ from
this task's fresh diagonalization matches the corresponding mesh-neighbour overlap task 9000
already exports for the same pair, checked for band 1 specifically -- not the full occupied
window, since bands 2--4 are triply degenerate at $\Gamma$ (Part~\ref{part:berry}) and two
\emph{independent} diagonalizations are free to pick different, equally valid bases within
that degenerate subspace; comparing raw overlap-matrix elements of a degenerate window
between two independent diagonalizations is not meaningful (confirmed empirically: bands 1--2
agree to the same $\sim 10^{-3}$ floor between the two methods, bands 3--4 do not, consistent
with an arbitrary within-subspace unitary mixing rather than a convention bug). Only
gauge-invariant quantities built from a whole degenerate window (e.g.\ Part~\ref{part:berry}'s
FHS flux) are safe to compare that way (\texttt{tests/test\_calculation\_eigenstates.py}).

\subsection{How to use in code}

\begin{verbatim}
with calc.eigenstate_session() as session:
    state = session.get_eigenstates((0.1, 0.2, 0.05))
    state.energies  # (nstsv,) Hartree
    state.evecsv    # (nstsv, nstsv) complex, evecsv[:, i] the i-th eigenvector

    m = session.overlap((0.0, 0.0, 0.0), (0.1, 0.0, 0.0), ist0=1, ist1=4)
    # m[a, b] = <psi_{1+a}(k_a)|psi_{1+b}(k_b)>
\end{verbatim}
\texttt{get\_eigenstates(k)}/\texttt{get\_overlap(k\_a, k\_b, ist0, ist1)} are one-off
convenience wrappers around a short-lived session, for a single query.

\part{Quantum geometry: the quantum metric, alongside Berry curvature}
\label{part:quantumgeometry}

\section{Setting: the quantum geometric tensor}

Berry curvature (Part~\ref{part:berry}, Eq.~\eqref{eq:berry}) is only half of the local
geometric information carried by a family of Bloch states $\psi(\mathbf k) =
(|n_1(\mathbf k)\rangle,\dots,|n_J(\mathbf k)\rangle)$ (a non-degenerate band, or a gapped
$J$-band manifold). The full object is the quantum geometric tensor (Provost \& Vallee,
\emph{Riemannian structure on manifolds of quantum states}, Commun.\ Math.\ Phys.\
\textbf{76}, 289 (1980)),
\begin{equation}
  Q_{ab}(\mathbf k) = \mathrm{Tr}\bigl[\partial_a P(\mathbf k)\, Q(\mathbf k)\,
    \partial_b P(\mathbf k)\bigr], \qquad
  P(\mathbf k) = \psi(\mathbf k)\,\psi^\dagger(\mathbf k), \quad Q = \mathbb 1 - P,
  \label{eq:qgt}
\end{equation}
the projector onto the $J$-dimensional manifold and its complement. $Q_{ab}$ decomposes
into a real symmetric part and an imaginary antisymmetric part,
\begin{equation}
  Q_{ab}(\mathbf k) = g_{ab}(\mathbf k) - \tfrac{i}{2} F_{ab}(\mathbf k), \qquad
  g_{ab} = \mathrm{Re}\,Q_{ab}, \quad F_{ab} = -2\,\mathrm{Im}\,Q_{ab},
  \label{eq:qgt-decomp}
\end{equation}
where $F_{ab}$ is exactly the (non-Abelian) Berry curvature of Eq.~\eqref{eq:berry}, and
$g_{ab}$ is the \emph{quantum metric} (Fubini-Study metric restricted to the manifold): a
genuinely distinct, gauge-invariant observable, not derivable from $F_{ab}$ alone. Where
$F_{ab}$ measures the manifold's holonomy (a phase acquired going around a loop),
$g_{ab}$ measures how much the manifold itself \emph{rotates} in Hilbert space between
$\mathbf k$ and $\mathbf k+d\mathbf k$ -- physically, it sets the quantum (Fubini-Study)
distance between neighbouring Bloch states, and appears directly in e.g.\ the
lower bound on the spread of maximally localized Wannier functions, the semiclassical
correction to the density of states in a strong field, and superfluid weight/quantum-metric
contributions to flat-band superconductivity.

A direct consequence of Eq.~\eqref{eq:qgt}, worth stating since it becomes a concrete
verification check below: $Q_{ab} = \sum_{n=1}^J \langle\partial_a n(\mathbf k)|\,Q(\mathbf
k)\,|\partial_b n(\mathbf k)\rangle$ (sum over the $J$ states spanning the manifold) is a
sum of Gram-matrix-like terms $\langle u_a|Q|u_b\rangle$, each individually positive
semi-definite as a $2\times2$ Hermitian matrix in the direction indices $(a,b)$ (since
$Q=Q^\dagger=Q^2$), and a sum of PSD matrices is PSD. A $2\times2$ Hermitian PSD matrix has
non-negative determinant, so with Eq.~\eqref{eq:qgt-decomp},
\begin{equation}
  \det Q(\mathbf k) = g_{11}g_{22} - g_{12}^2 - \bigl(F_{12}/2\bigr)^2 \ge 0
  \quad\Longrightarrow\quad
  \det g(\mathbf k) \ge \Bigl(\frac{F_{12}(\mathbf k)}{2}\Bigr)^{\!2},
  \label{eq:qgt-psd-bound}
\end{equation}
an exact inequality, not an order-of-magnitude estimate -- the metric can never be "too
small" for the curvature it accompanies at the same point. Both $g_{ab}$ and $F_{ab}$ are invariant
under $\psi(\mathbf k)\to\psi(\mathbf k)\,U(\mathbf k)$ for any $\mathbf k$-dependent
$U(J)$ gauge rotation of the manifold, since Eq.~\eqref{eq:qgt} is built entirely from the
gauge-invariant projector $P$, never from $\psi$ directly.

\section{Discretization from nearest-neighbour overlaps}

As in Part~\ref{part:berry}, $\psi(\mathbf k)$ is only ever known up to an arbitrary,
numerically-determined per-$\mathbf k$-point gauge, so Eq.~\eqref{eq:qgt} cannot be
evaluated by literally differentiating $P(\mathbf k)$. The discretized quantum metric
used here is the one already implicit in Marzari \& Vanderbilt's gauge-invariant
Wannier spread functional (PRB \textbf{56}, 12847 (1997), Eqs.~27, 30; reviewed in
Marzari, Mostofi, Yates, Souza \& Vanderbilt, Rev.\ Mod.\ Phys.\ \textbf{84}, 1419
(2012) \S II.C) -- the same overlap matrix
\begin{equation}
  M_{ab}(\mathbf k, \mathbf k+\mathbf v) = \langle n_a(\mathbf k) | n_b(\mathbf
  k+\mathbf v) \rangle
  \label{eq:qg-overlap}
\end{equation}
that Eq.~\eqref{eq:link}'s link variable is built from, but now entering through its
\emph{modulus} rather than its phase:
\begin{equation}
  D(\mathbf v) \equiv J - \mathrm{Re}\,\mathrm{Tr}\bigl[M(\mathbf k,\mathbf k+\mathbf v)\,
    M(\mathbf k,\mathbf k+\mathbf v)^\dagger\bigr]
  = J - \mathrm{Tr}\bigl[P(\mathbf k)\,P(\mathbf k+\mathbf v)\bigr].
  \label{eq:qdist}
\end{equation}
The second equality is Resta's projector form: $M M^\dagger$ has matrix elements
$\langle n_a(\mathbf k)|P(\mathbf k+\mathbf v)|n_{a'}(\mathbf k)\rangle$, so its trace is
$\mathrm{Tr}[P(\mathbf k)P(\mathbf k+\mathbf v)]$ directly, and
$D(\mathbf v) = \mathrm{Tr}[P(\mathbf k)\,Q(\mathbf k+\mathbf v)]$ -- the same bracket
Marzari-Vanderbilt sum over neighbour shells (weighted by their finite-difference
weights $w_{\mathbf v}$) to get the gauge-invariant part of the total spread, $\Omega_I
= \sum_{\mathbf k,\mathbf v} w_{\mathbf v}\, D(\mathbf v)$.

$D(\mathbf v)$ is exactly the quadratic form of the metric on $\mathbf v$: expanding
$P(\mathbf k+\mathbf v) = P + v^a\partial_a P + \tfrac12 v^av^b\partial_a\partial_b P +
O(v^3)$ and using idempotency ($P^2=P$, so $\mathrm{Tr}[P\partial_a P]=0$ and
$\mathrm{Tr}[P\,\partial_a\partial_b P] = -\mathrm{Tr}[\partial_a P\,\partial_b P]$, both
derived by differentiating $P^2=P$ once and twice and tracing) gives
\begin{equation}
  \mathrm{Tr}[P(\mathbf k)P(\mathbf k+\mathbf v)] = J - \tfrac12 v^a v^b\,
    \mathrm{Tr}[\partial_a P\,\partial_b P] + O(v^3),
\end{equation}
and a second idempotency identity, $\mathrm{Tr}[P\,\partial_aP\,\partial_bP] +
\mathrm{Tr}[P\,\partial_bP\,\partial_aP] = \mathrm{Tr}[\partial_aP\,\partial_bP]$
(differentiate $P\partial_aP+\partial_aP\,P=\partial_aP$ appropriately), shows the
symmetric combination $\tfrac12\mathrm{Tr}[\partial_aP\partial_bP]$ equals
$\mathrm{Tr}[P\,\partial_aP\,\partial_bP]$ symmetrized in $(a,b)$ -- i.e.\ exactly
$g_{ab}$ of Eq.~\eqref{eq:qgt-decomp} (using $Q\partial_bP = \partial_bP - P\partial_bP$
and idempotency once more to rewrite $\mathrm{Tr}[\partial_aP\,Q\,\partial_bP]$ in this
form). So, to leading order in $|\mathbf v|$,
\begin{equation}
  D(\mathbf v) = g_{ab}(\mathbf k)\, v^a v^b + O(|\mathbf v|^3).
  \label{eq:qdist-quadratic}
\end{equation}
Evaluating $D$ along the two loop directions $\mathbf v_1=\texttt{dk}\,\mathbf b_1$,
$\mathbf v_2=\texttt{dk}\,\mathbf b_2$ gives the diagonal components directly,
$g_{11}=D(\mathbf v_1)/|\mathbf v_1|^2$, $g_{22}=D(\mathbf v_2)/|\mathbf v_2|^2$; the
off-diagonal component follows from the polarization identity for the bilinear form
$D$,
\begin{equation}
  g_{12} = \frac{D(\mathbf v_1+\mathbf v_2) - D(\mathbf v_1) - D(\mathbf v_2)}
    {2\,|\mathbf v_1|\,|\mathbf v_2|},
  \label{eq:g12-polarization}
\end{equation}
needing one extra overlap, $M(\mathbf k,\mathbf k+\mathbf v_1+\mathbf v_2)$, beyond the
four already needed for $F_{12}$'s closed loop (Eq.~\eqref{eq:link}, corners $\mathbf k$,
$\mathbf k+\mathbf v_1$, $\mathbf k+\mathbf v_1+\mathbf v_2$, $\mathbf k+\mathbf v_2$,
anchored rather than centered, the same convention \texttt{parsers.berry}'s mesh mode --
not its centered path mode -- uses). Because $\det(S^{-1/2})$ for any Hermitian
positive-definite $S$ is real and positive (\S~below), this construction leaves $F_{12}$
computed from the same corners numerically identical to Part~\ref{part:berry}'s
Eq.~\eqref{eq:fieldstrength}.

\section{elkpy's implementation: Löwdin normalization, and why curvature didn't need it}

Elk's \texttt{genolpq} overlap (Part~\ref{part:eigenstates}, Eq.~\eqref{eq:overlap-integral})
carries a real-space truncation floor of order $10^{-3}$: $O(\mathbf k,\mathbf k)$ is the
identity only to that tolerance, not exactly (\S14 of \texttt{docs/design.md}). Write the
true overlap as $M = (1-\epsilon)M_{\rm exact}$ for a small common-mode deficiency
$\epsilon$. Part~\ref{part:berry}'s curvature is immune to this: $\arg\det[(1-\epsilon)
M_{\rm exact}] = \arg\det M_{\rm exact}$ exactly, since $(1-\epsilon)^{J}$ is a positive
real scalar and cannot shift an argument. Eq.~\eqref{eq:qdist} is not: $D(\mathbf v) =
J - (1-\epsilon)^2\,\mathrm{Tr}[M_{\rm exact}M_{\rm exact}^\dagger] \approx 2\epsilon J +
g_{ab}v^av^b$ to first order in $\epsilon$ -- a constant offset that does not shrink with
$\mathbf v$, so it \emph{dominates} the genuine $O(|\mathbf v|^2)$ metric signal as
\texttt{dk}$\,\to 0$, the opposite of a numerical artifact that simply averages out.
Measured directly (bulk Si, generic k-point): the raw (unnormalized) $g_{11}$ at
\texttt{dk} = 0.05, 0.02, 0.01, 0.005, 0.002 comes out 40.5, 48.2, 52.5, 62.9, 131.8 Bohr$^2$
-- diverging, the $1/\texttt{dk}^2$ signature of an uncorrected constant offset -- while
the value below (Löwdin-normalized) comes out 40.5, 47.6, 49.7, 50.5, 51.0, converging.

The fix is Löwdin symmetric normalization: query each loop corner's own self-overlap
$S(\mathbf c) = O(\mathbf c,\mathbf c)$ (already meaningful and available via
\texttt{EigenstateSession.overlap()} with $\mathbf k_a=\mathbf k_b$ -- the one new
ingredient curvature alone never needed a name for) and replace every cross overlap by
\begin{equation}
  \tilde M(\mathbf c_1,\mathbf c_2) = S(\mathbf c_1)^{-1/2}\, M(\mathbf c_1,\mathbf c_2)\,
    S(\mathbf c_2)^{-1/2}
  \label{eq:lowdin}
\end{equation}
before using it in Eq.~\eqref{eq:qdist} or Eq.~\eqref{eq:link}. This is exact by
construction, not a leading-order correction: the normalized self-overlap
$S^{-1/2}SS^{-1/2}=\mathbb 1$ identically at every corner, whatever $\epsilon$ actually
is at that corner (it need not even be the same at every corner, unlike the illustrative
common-mode $\epsilon$ above). It is inert for curvature: $S^{-1/2}$ is Hermitian
positive-definite, so $\det(S^{-1/2})$ is real and positive, and
$\arg\det\tilde M = \arg\det M$ exactly -- Eq.~\eqref{eq:lowdin} is applied uniformly to
both quantities in \texttt{parsers.quantum\_geometry.compute\_quantum\_geometry()}
purely for uniformity, not because curvature needs it.

No new Fortran task was needed for any of this (unlike Parts~\ref{part:berry}--
\ref{part:eigenstates}, each of which added one): every ingredient --
Eq.~\eqref{eq:qg-overlap}'s cross overlaps and now also each corner's self-overlap -- is
already exposed by the task 9002 interactive session (Part~\ref{part:eigenstates}), so
\texttt{Calculation.get\_quantum\_geometry()} drives everything above purely from
Python-side \texttt{EigenstateSession.overlap()} calls (9 per requested k-point: 4
self-overlaps, 5 cross overlaps) plus the discretization arithmetic in
\texttt{parsers/quantum\_geometry.py}, independently unit-testable against synthetic
data with no Elk run at all
(\texttt{tests/test\_quantum\_geometry\_gauge\_invariance.py}), the same way
\texttt{parsers/berry.py} is.

A remaining discretization subtlety worth stating plainly, not a bug: $g_{12}$
(Eq.~\eqref{eq:g12-polarization}) is a difference of three comparable-magnitude
$D$ values, so its own $O(\texttt{dk})$ discretization error is amplified relative to
its $O(\texttt{dk}^2)$ signal in a way $g_{11}$/$g_{22}$'s direct evaluation is not --
observed on monolayer h-BN's K/K$'$ valleys (\S~below): $g_{12}(K)$ and $g_{12}(K')$,
required equal in the $\texttt{dk}\to 0$ limit by time-reversal symmetry ($g_{ab}(-\mathbf
k)=g_{ab}(\mathbf k)$, the same argument Part~\ref{part:berry} uses for curvature's sign
flip, applied here to a quantity that does \emph{not} flip sign), only visibly converge
toward each other as \texttt{dk} is refined, unlike $g_{11}$/$g_{22}$, which already
agree to $<0.1\%$ at a single practical \texttt{dk}.

\subsection{Verified}

On synthetic data (no Elk run): $g_{ab}$, $F_{ab}$ and $Q_{ab}$ are invariant under a
random \emph{non-diagonal} $U(J)$ gauge mixing at every loop corner -- the non-Abelian
generalization of Part~\ref{part:berry}'s diagonal-only gauge check; that Löwdin
normalization exactly removes a synthetic uniform truncation deficiency applied to
otherwise-identical states (the true metric is exactly zero in that construction, and
only the normalized result recovers it); and, pinning the \emph{absolute} normalization
(not just invariance) against an analytically known case -- the spin coherent state
$|n(\theta,\phi)\rangle = (\cos\tfrac\theta2, e^{i\phi}\sin\tfrac\theta2)$, Provost \&
Vallee's own original worked example -- the metric converges to the exact round-sphere
result $g = \tfrac14\,\mathrm{diag}(1,\sin^2\theta)$ and the curvature to
$F_{\theta\phi}=\tfrac12\sin\theta$ (the sign this module's overlap-matrix convention
produces; magnitude is the solid-angle/2 result standard for a spin-1/2)
(\texttt{tests/test\_quantum\_geometry\_gauge\_invariance.py}).

Against a real compiled binary: the $\texttt{dk}$ divergence-vs-convergence contrast
above (bulk Si); that curvature from \texttt{get\_quantum\_geometry()} agrees with
\texttt{get\_berry\_curvature\_path()} (task 9001) on literally identical loop corners
(anchoring this method's loop at task 9001's own corner 1, with double the step size,
walks the same four points in the same cyclic order -- Part~\ref{part:berry}'s own
path/mesh cross-check applied one level up), confirming the new Python-level loop
construction introduces no sign/corner-order bug independent of the already-verified FHS
arithmetic it reuses; and, on monolayer h-BN's occupied valence manifold (same structure
as Part~\ref{part:berry}'s K/K$'$ verification): the quantum metric is positive
semi-definite at $\Gamma$, K, M, K$'$; curvature vanishes (to $<1\%$ of the valley scale)
at the time-reversal-invariant momenta $\Gamma$ and M and is K/K$'$-antisymmetric,
reproducing Part~\ref{part:berry}'s own finding via an entirely independent Python code
path (\texttt{EigenstateSession.overlap()} rather than task 9001's dedicated corners);
$g_{11}$/$g_{22}$ are K/K$'$-symmetric to $<1\%$ at $\texttt{dk}=0.01$; and, at all four
points, the exact bound Eq.~\eqref{eq:qgt-psd-bound} holds -- with comfortable margin at
K/K$'$ ($\det g$ about twice the bound) and trivially at $\Gamma$/M (both sides
near zero) -- tying the new metric's absolute scale to the already-trusted curvature
value at the same point without dividing by a curvature that vanishes at two of the four
points (\texttt{tests/test\_calculation\_quantum\_geometry.py}).

\subsection{How to use in code}

\begin{verbatim}
result = calc.get_quantum_geometry(
    [(0.0, 0.0, 0.0), (1 / 3, 1 / 3, 0.0)],  # Gamma, K -- fractional coordinates
    1, 4, directions=(1, 2), dk=0.01,
)
result[1]["g"]                # (2,2) real array [[g11,g12],[g12,g22]], Bohr^2
result[1]["berry_curvature"]  # Bohr^-2, identical convention to get_berry_curvature_path
result[1]["Q"]                # (2,2) complex array, Q = g - (i/2)*berry_curvature*[[0,1],[-1,0]]
\end{verbatim}
Same \texttt{kpoints=}/\texttt{kpath=}/\texttt{ist0}/\texttt{ist1}/\texttt{directions}/
\texttt{dk}/\texttt{npoints} conventions as \texttt{get\_berry\_curvature\_path()}
(Part~\ref{part:berry}), including symbolic \texttt{kpath="GKMK'G"}-style paths.

\part{Atom-projection operators}
\label{part:atomprojection}

\section{Setting: the muffin-tin partition of the unit cell}

An all-electron LAPW calculation partitions the unit cell into non-overlapping muffin-tin
(MT) spheres, one centered on each atom, plus the interstitial region between them
(Singh \& Nordstr\"om, \emph{Planewaves, Pseudopotentials, and the LAPW Method}, 2nd ed.,
Springer, 2006, ch.~2) -- the same partition Part~\ref{part:eigenstates}'s APW+lo basis is
built on. This partition is exact and exhaustive: every point in the cell lies in exactly
one region, so for any two states $\psi_i,\psi_j$ of the same diagonalization,
\begin{equation}
  \langle\psi_i|\psi_j\rangle
  = \sum_{\alpha=1}^{N_{\rm atom}} \int_{{\rm MT}_\alpha} \psi_i^*\psi_j\,d^3r
  \;+\; \int_{\rm interstitial} \psi_i^*\psi_j\,d^3r,
  \label{eq:mt-partition}
\end{equation}
with no double-counted or omitted volume. Define the atom-projection operator $P_\alpha$
as the restriction of the identity to atom $\alpha$'s sphere,
\begin{equation}
  (P_\alpha)_{ij} = \langle\psi_i|\hat P_\alpha|\psi_j\rangle
  = \int_{{\rm MT}_\alpha} \psi_i^*(\mathbf r)\,\psi_j(\mathbf r)\,d^3r,
  \label{eq:atomproj}
\end{equation}
an $n_{\rm st}\times n_{\rm st}$ matrix in whatever band window $\{\psi_i\}$ is requested.
$P_\alpha$ is Hermitian by construction (Eq.~\eqref{eq:atomproj} is a restricted inner
product) and positive semi-definite (it is literally a Gram matrix of the states
restricted to a sub-volume -- see the discretized form below, where positivity is
manifest). Eq.~\eqref{eq:mt-partition} restated in this notation is the operator identity
\begin{equation}
  \sum_{\alpha=1}^{N_{\rm atom}} P_\alpha + P_{\rm interstitial} = \mathbb 1,
  \label{eq:atomproj-partition}
\end{equation}
the exact-partition constraint this Part's verification below checks directly: summing
every atom's matrix can only ever fall short of the identity (never exceed it), by
exactly the interstitial region's own weight.

Diagonal entries $(P_\alpha)_{ii}\in[0,1]$ are the familiar atom-projected/partial weight
of state $i$ on atom $\alpha$ -- the quantity Elk's own atom-and-$\ell$-resolved partial
DOS (\texttt{dos} task, \texttt{PDOS\_Sss\_Aaaaa.OUT}) and band-character output
(\texttt{bandstr} tasks 21--24, \texttt{BAND\_Sss\_Aaaaa.OUT}) already compute and
energy-bin or tie to a band path. Off-diagonal entries $(P_\alpha)_{ij}$, $i\neq j$, are
the genuinely new piece: matrix elements of an actual operator between two \emph{different}
eigenstates of the same diagonalization, not merely a per-state weight -- e.g.\ they let
$P_\alpha$ be applied to an arbitrary superposition $c=\sum_i c_i\,|\psi_i\rangle$
(expressed as a coefficient vector in this same band window) via ordinary matrix-vector
multiplication, $P_\alpha c$, rather than only to single basis states.

\section{Discretized form: reusing \texttt{wfmtsv} and \texttt{wr2cmt}}

Elk already expands a batch of second-variational states into atom $\alpha$'s
muffin-tin, resolved by $(\ell m,\sigma)$ and radial shell, via upstream
\texttt{wfmtsv.f90} -- exactly the ingredient \texttt{gendmatk.f90} uses to build the
atom/$(\ell,m)$-resolved density-matrix entries behind \texttt{dos}/\texttt{bandstr}'s
output, and the ingredient this Part reuses unchanged. Writing the expansion as
$\psi_{i,\ell m\sigma}(r)$ (radial function on shell index, angular/spin index packed
separately) and $w(r)$ for the \texttt{wr2cmt} quadrature weight (which depends only on
radial shell, not on $\ell,m,\sigma$),
\begin{equation}
  (P_\alpha)_{ij} = \sum_{\ell m,\sigma}\int_0^{R_\alpha}
    \psi_{i,\ell m\sigma}^*(r)\,\psi_{j,\ell m\sigma}(r)\,w(r)\,dr,
  \label{eq:atomproj-discrete}
\end{equation}
the muffin-tin-sphere restriction of $\langle\psi_i|\psi_j\rangle$ summed over every
$(\ell,m,\sigma)$ channel (the total atomic weight, not an $\ell$-resolved breakdown --
\texttt{gendmatk} keeps the $(\ell,m)$ index open since \texttt{dos}/\texttt{bandstr}
want an $\ell$-resolved curve; nothing here needs that, so it is summed away).
Eq.~\eqref{eq:atomproj-discrete}, with $w(r)>0$ a genuine quadrature weight, is manifestly
a weighted Gram matrix -- Hermitian and PSD without further argument, since it is
$\Psi^\dagger W \Psi$ for $\Psi$ the (radial-shell,$\ell m\sigma$)$\times n_{\rm st}$
coefficient matrix and $W={\rm diag}(w)\succeq0$.

\texttt{gendmatk} only ever evaluates the \emph{diagonal} $i=j$ (one state's own
occupation-matrix entry, needed for LDA+U-style applications), and only ever behind a
task (\texttt{dos}, \texttt{bandstr}) that reduces the result before it reaches the
caller -- energy-binned into a DOS curve, or tied to Elk's own band-path generator, never
exposed raw at an arbitrary $\mathbf k$-point. \texttt{elkpy\_atomproj}
(\texttt{patches/0004-atom-projection.patch}, new subroutine in
\texttt{src/elkpy\_eigenstates.f90}) reuses \texttt{wfmtsv} identically -- no new radial
function, matching coefficient, or numerical method -- but evaluates the full
$i,j$ off-diagonal sum of Eq.~\eqref{eq:atomproj-discrete} as a single \texttt{zgemm} per
spin channel, at a fresh on-the-fly diagonalization (the same
\texttt{eveqnfv}/\texttt{eveqnsv} pattern as \texttt{elkpy\_diagonalize}/
\texttt{elkpy\_wfcorner}, Part~\ref{part:eigenstates}) at an arbitrary requested
$\mathbf k$-point.

\section{Why every atom in one query}

As with \texttt{evecsv} generally (Part~\ref{part:eigenstates}), a matrix from one
diagonalization is not comparable to a matrix from a \emph{separate} diagonalization,
even nominally at the same $\mathbf k$ -- nothing guarantees two independent
diagonalizations picked the same internal basis. Eq.~\eqref{eq:atomproj-partition} is
therefore only checkable when every $P_\alpha$ being summed came from the same
diagonalization. The task 9002 \texttt{PROJECTION $k_1\,k_2\,k_3\,ist_0\,ist_1$} query
returns all $N_{\rm atom}$ matrices from one diagonalization, looping
\texttt{elkpy\_atomproj} over every atom server-side (reusing that one diagonalization's
\texttt{apwalm}/\texttt{evecfv}/\texttt{evecsv} for each atom's \texttt{wfmtsv} call --
one extra \texttt{wfmtsv} call per atom, not another diagonalization) rather than
letting a caller assemble them from separate one-atom-at-a-time queries.

\subsection{Verified}

Against a real compiled binary (bulk Si, generic $\mathbf k$-point): every returned
$P_\alpha$ is Hermitian and positive semi-definite, a direct algebraic consequence of
Eq.~\eqref{eq:atomproj-discrete}'s Gram-matrix form rather than external reference data;
$\mathbb 1-\sum_\alpha P_\alpha$ (Eq.~\eqref{eq:atomproj-partition}'s interstitial
remainder) is itself Hermitian PSD, i.e.\ the sum never overshoots the identity, with
each atom's own diagonal weight a substantial, physically reasonable fraction of the
cell (neither near zero nor exceeding 1) for a tetrahedrally-bonded solid like Si; and
diamond Si's two atoms -- related by inversion through the bond midpoint, which sends
$\mathbf k\to-\mathbf k$ and therefore fixes $\Gamma$ -- have identical weight on band 1
at $\Gamma$ (non-degenerate; bands 2--4 are degenerate there, the same caveat
Part~\ref{part:eigenstates} documents for comparing eigenstates of a degenerate
subspace). A diagonal entry agrees, to 5 decimal places, with an entirely independent
Fortran code path -- upstream \texttt{bandstr.f90} task 21's own
\texttt{BAND\_Sss\_Aaaaa.OUT} "total atomic character" column, which calls the same
\texttt{gendmatk}/\texttt{wfmtsv} machinery via a separate call site (Elk's own path
diagonalization, not \texttt{elkpy\_atomproj}'s fresh one) and a separate reduction (an
explicit per-$\ell$ loop, not this feature's \texttt{zgemm}), once task 21's default
$\ell_{\rm max}^{\rm db}=3$ is raised to match the ground state's own $\ell_{\rm
max}^{\rm o}=6$ -- this catches what the Hermitian/PSD and sum-below-identity checks
alone cannot: both are satisfied by any internally consistent normalization, even one
that silently undercounts weight (e.g.\ a packing/stride bug in the \texttt{zgemm}
reduction), so only an external reference computing the same number a different way
closes that gap. A spin-polarized run ($n_{\rm spinor}=2$) remains Hermitian PSD and
sum-below-identity, exercising the spin-channel accumulation the unpolarized checks
above never run twice (spin-orbit coupling, Part~\ref{part:soc}, remains untested for
this feature). On monolayer h-BN (same structure as Part~\ref{part:berry}'s K/K$'$
verification), at $K=(1/3,1/3,0)$ the occupied valence-top ($\pi$) band is
N-dominated and the unoccupied conduction-bottom ($\pi^*$) band is B-dominated -- the
standard qualitative picture for h-BN's band character (the more electronegative N
pulling the bonding state's weight toward itself), a sharp sign-of-the-effect
prediction rather than a plausibility band, the same spirit as Part~\ref{part:berry}'s
K/K$'$ curvature-antisymmetry check
(\texttt{tests/test\_calculation\_atom\_projection.py}).

\subsection{How to use in code}

\begin{verbatim}
proj = calc.get_atom_projection((0.1, 0.2, 0.05), ist0=1, ist1=4)
proj.matrices.shape          # (natmtot, nst, nst) complex
proj.matrices[a][i, i].real  # state (ist0+i)'s fractional weight on atom a's muffin tin

n = calc.structure.atom_index("N")   # (species, index) -> this array's atom axis
proj.matrices[n]                     # atom-projection operator for that N atom
\end{verbatim}
\texttt{EigenstateSession.atom\_projection(k, ist0, ist1)} is the underlying session
method (\texttt{eigenstate\_session()}, Part~\ref{part:eigenstates}); this is a one-off
convenience wrapper opening and closing its own session, for a single query.

\begin{thebibliography}{9}
\bibitem{koelling1977}
D.~D.~Koelling and B.~N.~Harmon,
\emph{A technique for relativistic spin-polarised calculations},
J.\ Phys.\ C: Solid State Phys.\ \textbf{10}, 3107 (1977).
\bibitem{singhnordstrom2006}
D.~J.~Singh and L.~Nordstr\"om,
\emph{Planewaves, Pseudopotentials, and the LAPW Method}, 2nd ed.,
Springer (2006).
\end{thebibliography}

\end{document}
